\documentclass[11pt,a4paper]{report}

\usepackage{xcolor}
\def\farbe{cyan}

\usepackage{dclecture}

\usepackage{pigpen}
\pdfgentounicode=0

\usepackage{morse}
\def\lm{\Large\morse}

\setlength{\headheight}{24pt}



%%% Fancy Header and Footer
\renewcommand{\headrule}{\vbox to 0pt{\hbox to\headwidth{\color{\farbe}\rule{\headwidth}{1pt}}\vss}}
\pagestyle{fancy} %eigener Seitenstil
\fancyhf{} %alle Kopf- und Fusszeilenfelder bereinigen
\fancyhead[C]{Topic} %Kopfzeile mitte
%\fancyhead[R]{\includegraphics[width=0.2cm]{x.png}}
\fancyfoot[C]{\thepage}







\begin{document}

\section*{Murder in the Mathematics Department}
A murder has been committed in the Maths department at a
nearby school.

Your job is to decode the clues to find
\begin{enumerate}
    \item the identity of the murderer
    \item  the murder weapon
   \item the room in which the murder took place

\end{enumerate}

The seven suspects are:
\begin{multicols}{4}
    \begin{itemize}
        \item Mr Broome 
        \item Dr France 
        \item Mr Halai 
        \item Miss Kaur
        \item Mr Scotland 
        \item Mr Sheppard 
        \item Mr Smith
    \end{itemize}
\end{multicols}
The possible murder weapons are:
\begin{multicols}{4}
    \begin{itemize}
        \item Metre ruler 
        \item Chair 
        \item Scissors
        \item Compass 
        \item Textbook 
        \item Pencil
        \item Stapler
    \end{itemize}
\end{multicols}
The room in which the murder was committed could be:
\begin{verbatim}
    16 23 24 25 28 29 Computer room
\end{verbatim}



\newpage

\section*{Code 1: Pigpen}
The first code uses the following key. You need to work out how to use the key
to decode the message which follows.

\centfig{0.8}{pigpen.png}

{\pigpenfont THE ROOM IN WHICH THE MURDER WAS COMMITTED HAS A ROOM NUMBER}

\section*{Code 2: Polybius Square}
\centfig{0.5}{polybius}

\begin{verbatim}
    (5,2)(3,4)(5,5) (3,3)(1,1)(3,2)(4,5)(5,5)(3,2)(5,5)(3,2) 
    (4,5)(5,3)(5,5)(4,2) (4,3)(5,3)(5,2) (3,4)(1,5)(2,1)(5,5) 
    (1,5)(4,3) (5,5) (4,4)(4,3) (5,2)(3,4)(5,5)(4,4)(3,2) 
    (4,3)(1,5)(3,3)(5,5)
\end{verbatim}

\section*{Code 3}
\begin{verbatim}
    20,8,5  13,21,18,4,5,18  23,1,19  14,15,20  9,14 1 
    16,18,9,13,5  14,21,13,2,5,18,5,4  18,15,15,13
\end{verbatim}

\section*{Code 4}
Your only clue for this one is that it was received by mobile phone.

\centfig{0.4}{mobile}

\begin{verbatim}
    843068733709327660363706680428302074277063825072780
\end{verbatim}

\section*{Code 5: Atbash}
Hint: every letter represents some other letter in a simple system

\begin{verbatim}
    GSV ILLN MFNYVI RH Z NFOGRKOV LU ULFI
\end{verbatim}

\section*{Code 6: Caesar Cipher}

\begin{verbatim}
    KYV DLIUVIVIJ ERDV NZCC KVCC PFL NYRK TFLEKIP YV ZJ WIFD
\end{verbatim}

\section*{Code 7: Morse Code}
\centfig{0.8}{morse}

{\lm the room number}

{\lm has eight factors}

\section*{Code 8: The Final Challenge}

{\lm kdc qxf mrm qn mx } 

{\lm rc ynaqjyb frcq }

{\lm bxvncqrwp cqjc}

{\lm  bcdmnwcb ljw brc xw}



\end{document}